\documentclass[UTF8,11pt,a4paper]{ctexart}
\newcommand{\DocTitle}{评分与提交}
\newcommand{\DocSubtitle}{公式、功能门槛与交付清单}
% Shared layout for the CW2 companion documents.
\usepackage[margin=2.1cm,headheight=16pt]{geometry}
\usepackage{amsmath,amssymb,booktabs,array,tabularx,ragged2e}
\usepackage{enumitem,xcolor,fancyhdr,listings,tcolorbox}
\tcbuselibrary{breakable}
\usepackage{hyperref}
\usepackage{bookmark}

\newcommand{\DocVersion}{3.0}
\newcommand{\EngineVersion}{1.1}
\newcommand{\DocDate}{2026-11-23}
\definecolor{ink}{HTML}{1F2937}
\definecolor{accent}{HTML}{92400E}
\definecolor{muted}{HTML}{4B5563}
\definecolor{linkink}{HTML}{174678}
\definecolor{codebg}{HTML}{F5F7FA}
\hypersetup{unicode,colorlinks=true,linkcolor=linkink,urlcolor=linkink,
  filecolor=linkink,pdftitle={SOE104 CW2 \DocTitle\space v\DocVersion},
  pdfauthor={Group A--E}}
\ctexset{section={format={\Large\bfseries\color{accent}}},
  subsection={format={\large\bfseries\color{ink}}}}
\pagestyle{fancy}
\fancyhf{}
\fancyhead[L]{\small\color{muted}SOE104 CW2 · Oregon Trail}
\fancyhead[R]{\small\color{muted}\DocTitle\ v\DocVersion}
\fancyfoot[C]{\small\color{muted}\thepage}
\renewcommand{\headrulewidth}{0.4pt}
\setlength{\parskip}{0.35em}
\setlength{\emergencystretch}{1.5em}
\renewcommand{\arraystretch}{1.23}
\setlist{itemsep=3pt,topsep=3pt,parsep=0pt,leftmargin=1.7em}
\newcolumntype{Y}{>{\RaggedRight\arraybackslash}X}
\newcolumntype{L}[1]{>{\RaggedRight\arraybackslash}p{#1}}
\newtcolorbox{keybox}[1]{breakable,title={#1},colback=orange!4,
  colframe=accent,boxrule=0.55pt,arc=2pt,fonttitle=\bfseries}
\newtcolorbox{warnbox}[1]{breakable,title={#1},colback=red!3,
  colframe=red!45!black,boxrule=0.55pt,arc=2pt,fonttitle=\bfseries}
\lstset{basicstyle=\small\ttfamily,backgroundcolor=\color{codebg},
  xleftmargin=0.9em,frame=l,rulecolor=\color{linkink!45},
  breaklines=true,columns=fullflexible,keepspaces=true,showstringspaces=false,
  aboveskip=0.65em,belowskip=0.65em}
\newcommand{\code}[1]{\texttt{#1}}
\newcommand{\source}[1]{\textup{\small\textcolor{muted}{[#1]}}}
\newcommand{\docref}[2]{\href{./#1.pdf}{#2}}
\renewcommand{\maketitle}{%
  \begin{center}
    {\LARGE\bfseries CW2 \DocTitle\par}
    \vspace{0.35em}{\large\DocSubtitle\par}
    \vspace{0.75em}{\small\color{muted}Group A--E · 整理：A
      \quad\DocDate\quad 文档 v\DocVersion\par}
    \vspace{0.2em}{\small\color{muted}Engine / Marking v\EngineVersion\par}
  \end{center}
  \vspace{0.4em}}


\begin{document}
\maketitle
\thispagestyle{fancy}

\begin{keybox}{评分先看运行是否有效}
本地用 JDK 21 验证编译，平台运行三个种子并给出回放和成绩。
常规分数按四项加权。编译失败、崩溃、超时和禁用 API 有单独的评分后果。
\end{keybox}

\section{总分和 Game 分}

\[
\mathrm{FINAL}=0.20\,\mathrm{OOP}+0.50\,\mathrm{Game}
  +0.15\,\mathrm{Naming}+0.15\,\mathrm{Efficiency}.
\]
Game 按三个种子的结果取平均。Marking 1.1 保留四项权重，更新了 Game。
记 $D$ 为终局里程，$N$ 为初始人数，$H$ 为初始成员的终局总 HP（死者计 0），
$U$ 为食品、牛、弹药和零件的终局总数量。
\[
\mathrm{Game}=40\frac{D}{2000}+30\frac{H}{100N}+30g(U).
\]
公告写明 $g(U)=U/1000$。平台对 $U\ge1000$ 是否另设封顶仍待验证。
下表使用公告公式和评分权重。 \source{TU} \source{TASK}
\begin{center}\small
\begin{tabularx}{\linewidth}{@{}L{2.2cm}L{4.6cm}Y@{}}
\toprule
贡献 & Game 算法 & 折算到 FINAL 的例子 \\
\midrule
到达 & 40 倍里程比例 & 走 2,000 英里得 20 分。 \\
状态 & 30 倍归一化平均 HP & 4 人初始、3 人满血，$H/(100N)=0.75$，得 11.25 分。 \\
货物 & 30 倍货物比例 & 600 单位实物得 9 分。现金不计入。 \\
\bottomrule
\end{tabularx}
\end{center}
例如抵达终点、归一化平均 HP 为 $0.80$、实物剩 600，Game 为 $40+24+18=82$。
若其他三项加权合计 50，FINAL 为 91。四人队少一个满血终局成员，状态分上限
比四人满血少 3.75 FINAL 分。实际受伤损失和后续行动价值另见\docref{04_decision_model}{《04 决策模型》}。

\section{代码质量与效率}

\begin{center}\small
\begin{tabularx}{\linewidth}{@{}L{3.0cm}Y@{}}
\toprule
维度 & 自查内容 \\
\midrule
OOP（20 分） & 子项为封装 25、继承 25、多态 25、接口 15、内聚 10。
字段受控、角色行为在各自类中实现，避免向下转型和类型判断链。 \\
Naming（15 分） & 类名用 PascalCase，方法与变量用 camelCase，名称体现职责。 \\
Efficiency（15 分） & 规则说明的平均每种子操作数门槛为
$\le8{,}000 \Rightarrow100$、$\le25{,}000 \Rightarrow80$、
$\le60{,}000 \Rightarrow60$、更高 $\Rightarrow$ 至多 40（效率子项分）。 \\
\bottomrule
\end{tabularx}
\end{center}
子项清单和实际反馈以课程平台为准。单次回调预算与 Efficiency 分档是两个量。
\source{ER} \source{TASK}

\clearpage
\section{功能门槛与沙箱}

\begin{center}\small
\begin{tabularx}{\linewidth}{@{}L{5.25cm}Y@{}}
\toprule
触发条件 & 后果或处理 \\
\midrule
编译失败、崩溃或超时 & 功能门槛生效，Game 为 0，FINAL 封顶 40。慎用可能抛异常的返回对象。 \\
单次回调超出约 100,000 ops & 超预算调用被跳过，回放会记录。单种子的墙钟上限为 10 秒。 \\
命中禁用 API & 提交被拒，得 0 分。引擎内部类不在学生 classpath 中。 \\
\bottomrule
\end{tabularx}
\end{center}
\subsection{源码审查时具体看什么}

\begin{center}\small
\begin{tabularx}{\linewidth}{@{}L{3.45cm}Y@{}}
\toprule
类别 & 用于审查的例子 \\
\midrule
反射与类加载 &
\code{java.lang.reflect.*}、\code{Class.forName}、\code{java.lang.invoke.*} 和自定义 \code{ClassLoader}。 \\
进程、退出和并发 &
进程和退出相关的 \code{Runtime}、\code{ProcessBuilder}、\code{System.exit}。
\newline 启动线程和使用 \code{java.util.concurrent.*} 也在禁用范围内。 \\
系统环境与 I/O &
\code{System.getenv}、\code{System.getProperty}、文件和网络访问，例如
\code{java.io.*}、\code{java.nio.*}、\code{java.net.*}。题面明示的 I/O 例外为 \code{System.out} 控制台输出。 \\
随机数与引擎状态 &
\code{java.util.Random} 禁用。\code{Math.random()} 与时间 API 由平台确定化。
行为通过返回对象交给引擎结算，禁止直接修改成员健康、物资和名单。 \\
\bottomrule
\end{tabularx}
\end{center}

\begin{keybox}{区分三个经常混在一起的结果}
队伍正常结束但没到终点时，按走过的距离与终局状态计分。
编译失败、崩溃或超时触发功能门槛。例如得不到有效 \code{Action} 类型时，
\code{new Action(null)} 会抛异常。命中禁用 API 则会被拒并计零分。
本地编译通过后，仍需看平台回放和反馈。 \source{TASK} \source{ER}
\end{keybox}

\clearpage

\section{提交清单}

\begin{enumerate}
  \item 在 LM 指定的评分平台提交。一队一份，使用本队的账户。
  \item 交 1--20 个 UTF-8 \code{.java} 文件，或把源文件放在单个 ZIP 的根目录。
  Java 使用默认包，每文件一个 public class，恰好一个 \code{PartyFactory} 实现，
  返回 3--5 名成员。适用的上传大小限制见平台，规则说明列出 20 MB。
  \item 课程截止为 \textbf{2026-12-13 23:59（中国时间）}。截止前可不限次提交，
  \textbf{只计最新一次}。题面列出迟交每日扣 5\%、累计至多 25\%，
  具体以平台说明为准。
  \item 课程要求提交 Java 源码，无书面报告。每名组员另在 LM 提交同伴评价，
  个人分可能按贡献调整。 \source{TASK} \source{ER}
\end{enumerate}

\noindent 打包时先打开 ZIP 核对层级，例如 \code{TeamFactory.java} 应在压缩包根目录，
无需先打开 \code{src/} 子目录。如果保留了课程示例的另一份 \code{PartyFactory}，
打包本组版本前要把示例排除。记录提交对应的 Git 版本，确认三个种子都有可核对的回放。

\section{学术诚信和 AI 辅助}

\begin{warnbox}{作者写代码，并能说明设计依据}
题面允许用 AI 讨论规则与 Java 概念、索取解释、提示或反馈。
实际提交的代码由组员自行编写并理解，每次使用 AI 辅助都按题面在代码文件头如实记录工具和用途。
课程禁止用 AI 代写作业提交。答辩与申报分别有各自的要求。
\end{warnbox}

平台做跨当前与历史队伍的代码相似度检查，并根据提交模式及代码相似度筛查可能由 AI 完成的作业。
被标记的队伍可能受邀做设计答辩，解释不了自己的提交将按学校学术诚信政策处理。
团队复核时，可让作者结合一个正常用例和一个边界例解释设计，提交记录只用于追溯改动。
\source{TASK}

\noindent 若用 AI 辅助，按实际情况在代码文件头记录。例如使用某工具理解接口说明时，
可以写下工具名和所问的具体主题。把括号内的内容换成真实信息。
\begin{lstlisting}
/* AI assistance: [tool name], used to clarify
   [the API rule or Java concept discussed]. */
\end{lstlisting}

\medskip
\noindent 资料口径：\source{TU} Trail Update 1.1；\source{TASK} CW2 Tasksheet
（交付、沙箱、评分和学术诚信）；\source{ER} Engine Rules Spec 第 9--10 节。

\end{document}
